% =====================================================================
%  Phase C --- Cryptography: what is signed, what is sealed, and why
% =====================================================================
\section{Authenticity and Confidentiality}
\label{sec:phaseC_security}

\subsection{Two different questions}

Habit says encrypt everything. Resisting that habit is what makes this
part of the design legible, so the two directions are treated
separately and the reasoning is written down.

\paragraph{Cloud $\rightarrow$ robot: signed, not encrypted} A request
says ``measure \station{P2,5R}''. That is not a secret. But issuing one
drives a robot around a greenhouse, and anyone who can reach the broker
can publish on the topic. The property that matters is therefore
\emph{authenticity}, not confidentiality. Requests carry an ECDSA-P256
signature over the canonical JSON; the robot refuses an unsigned request
and refuses one signed by the wrong key.

\paragraph{Robot $\rightarrow$ Cloud: sealed and signed} A report
contains the agronomic state of a commercial crop and a photograph of
it. It needs confidentiality \emph{and} origin, and gets both.

\subsection{The envelope}

Reports are sealed with ECIES over SECP256R1: an ephemeral key pair per
message, ECDH against the Cloud's public key, HKDF-SHA256 to derive the
content key, and AES-256-GCM for the payload. The ciphertext is then
signed with the robot's long-term ECDSA-P256 key. The complete wire
message is:

\begin{lstlisting}[style=bash]
{"schema":"agri-cloud/v1", "kind":"sealed",
 "robot":"youbot-01", "request_id":"...",
 "sent_at":"2026-08-08T19:23:27Z",
 "sealed": {
   "alg":  "ECIES-P256-HKDF-SHA256-AES256GCM",
   "epk":  "-----BEGIN PUBLIC KEY-----...",
   "nonce":"UgoQGAyNp7LauuRB",
   "ciphertext":"5k2ZWwLpBQVmV61fDxBg39Xy...",
   "sig":  "MEUCIGmZLjCjPHsyhdx97dGKKhmk..." }}
\end{lstlisting}

The routing metadata --- who sent it, for which request, when --- stays
in clear, because the broker and the operator need it to route and to
diagnose, and none of it is the measurement. The readings, the pose, the
QR image and the photograph are all inside the ciphertext.

\subsection{Sign-then-encrypt, or encrypt-then-sign}

The signature is over the \emph{ciphertext}. This is deliberate: it lets
the Cloud reject a forged message without first performing a private-key
operation on attacker-chosen input, and it means a tampered ciphertext
is detected by the cheap check rather than by AES-GCM's authentication
tag alone.

The suite pins six properties of this construction, and four of them are
refusals:

\begin{table}[!t]
\caption{Cryptographic properties asserted on every run. The refusals
matter more than the round trip: a system that accepts everything also
round-trips perfectly.}
\label{tab:crypto}
\centering\footnotesize
\begin{tabular}{@{}p{0.50\columnwidth}p{0.40\columnwidth}@{}}
\toprule
property & verdict \\
\midrule
a sealed payload round-trips              & \ok{accepted} \\
the plaintext is absent from the envelope & \ok{confirmed} \\
the wrong recipient key cannot open it    & \bad{refused} \\
a forged sender is rejected               & \bad{refused} \\
stripping the signature does not bypass the check & \bad{refused} \\
one flipped bit is detected               & \bad{refused} \\
the same plaintext seals differently every time & \ok{confirmed} \\
a rewritten target breaks the request signature & \bad{refused} \\
\bottomrule
\end{tabular}
\end{table}

The third-from-last matters more than it looks: an ephemeral key per
message means two identical readings produce two unrelated ciphertexts,
so an observer cannot tell that a station was measured twice with the
same result.

\subsection{Key custody, and the refusal that protects it}

The Cloud is meant to run on a different machine from the robot. That
intent is enforced rather than documented:

\begin{itemize}[leftmargin=*,itemsep=1pt,topsep=2pt]
  \item An \textbf{empty} key directory is bootstrapped with both pairs,
        so a first run on one machine works with no ceremony.
  \item A \textbf{populated} key directory is never bootstrapped behind
        the operator's back.
  \item \textbf{Neither side ever generates the other side's key.} The
        Cloud will not invent a robot identity; the robot will not
        invent a Cloud identity.
\end{itemize}

That last rule is the one that earns its place. Without it, a Cloud
started on a fresh machine with a forgotten key copy would generate a
robot public key of its own, and then reject every report from the real
robot --- with a valid signature failure, forever, and no indication
that the cause was a missing file. Instead it refuses to start and names
the file:

\begin{lstlisting}[style=bash]
/path/forgot/cloud_public.pem is missing.
\end{lstlisting}

The suite exercises the whole two-machine story: bootstrap an empty
directory, copy only the public keys to a second directory, start a
robot there, and assert that it runs \emph{and} that it does not hold
the Cloud's private key.

\subsection{Refusals are counted, not dropped}

A report that fails to open is not discarded silently. It is counted,
its reason is kept, and the dashboard displays it under
\textsc{rejected}. A pipeline that hides its rejections is a pipeline
that will one day be receiving nothing and look perfectly healthy.

In the reference campaign the panel read \emph{``nothing refused ---
every report verified''} across all \measured{48} reports, which is the
statement worth making: not that the cryptography was exercised, but
that it was exercised and nothing failed.
